\section{Three-clause refinements}\label{sec:three-clause}

For three literals per clause, a first moment at density six improves
the center bound enough to use the comparison interval \([0,8]\),
including for the main cap \(16n\). We also compare that statistic with
last-satisfiable and first-unsatisfiable indices capped at \(10n\).
Changing the cap of the last-satisfiable index has an exponentially
small effect on its normalized mean. Comparison with the first
unsatisfiable index also includes a one-step shift.

Throughout this section the clauses are independent uniformly signed
proper triples, sampled with replacement, and \(n\ge3\). Recall the
first unsatisfiable index \(H_n\), and write
\begin{equation}\label{eq:three-statistics}
 W_n=\min\{H_n-1,10n\},\qquad
 T_n=\min\{H_n,10n\},\qquad
 V_n=\min\{H_n-1,16n\}.
\end{equation}
Thus \(V_n\) is the statistic of the main proof at \(k=3\). Put
\(\nu_n=\EE W_n/n\), while \(\mu_n=\EE V_n/n\) as before.
For integer \(m\),
\begin{equation}\label{eq:three-survival}
\begin{aligned}
 P_n(m)&=\PP(W_n\ge m)&&\quad(0\le m\le10n),\\
 P_n(m)&=\PP(T_n>m)&&\quad(0\le m<10n).
\end{aligned}
\end{equation}
Only the last-satisfiable identity includes its capped endpoint.

\subsection{The density-six bound and changes of cap}

Each fixed assignment satisfies a proper triple with probability \(7/8\),
so \(P_n(m)\le2^n(7/8)^m\). Set
\[
 \sigma=2(7/8)^6<1,\qquad \rho=2(7/8)^{10}<1.
\]
The survival-sum identity for the positive integer \(H_n\) gives the
cap-independent bound
\begin{equation}\label{eq:three-h-mean}
 \EE H_n=\sum_{m\ge0}P_n(m)
 \le6n+\sum_{m\ge6n}2^n(7/8)^m
 =6n+8\sigma^n.
\end{equation}
In particular, the normalized means of \(W_n,T_n,V_n\) are all at most
seven for large \(n\). Splitting at \(6n\) also gives the direct estimate
\begin{equation}\label{eq:three-center}
 \nu_n\le6+4P_n(6n)\le6+4\sigma^n.
\end{equation}
The matching argument in Lemma~\ref{lem:center-bounds} gives
\(P_n(\lfloor c_0n\rfloor)=1-o(1)\) whenever \(c_0>0\) and
\(e^2c_0<1\). These clause counts lie below both caps, so there is
\(a_0>0\) such that
\begin{equation}\label{eq:three-anchors}
 a_0\le\nu_n\le\mu_n\le7
 \qquad\text{for all sufficiently large }n.
\end{equation}
The improved upper bound therefore applies to either cap.

The exact relations explain the distinction between first and last
indices at the capped endpoint. With the common cap \(10n\),
\begin{equation}\label{eq:three-same-cap}
 W_n=T_n-1+\ind_{\{H_n>10n\}}
     =T_n-1+\ind_{\{F_{n,10n}\in\sat\}}.
\end{equation}
If \(H_n\le10n\), the indices differ by one; otherwise both equal
\(10n\). To compare \(V_n\) with \(T_n-1\), count the extra satisfiable
prefixes:
\begin{equation}\label{eq:three-cross-cap}
\begin{gathered}
 V_n=T_n-1+E_n,\qquad
 E_n=\sum_{m=10n}^{16n}\ind_{\{H_n>m\}},\\
 0\le E_n\le(6n+1)\ind_{\{H_n>10n\}}.
\end{gathered}
\end{equation}
The first term is necessary: \(H_n=10n+1\) gives \(E_n=1\).
Since \(P_n(10n)\le\rho^n\), these identities imply
\begin{align}
 0\le\EE V_n-(\EE T_n-1)&\le(6n+1)\rho^n,
       \label{eq:three-mean-bridge}\\
 |\sd(V_n)-\sd(T_n)|&\le(6n+1)\rho^{n/2}.
       \label{eq:three-sd-bridge}
\end{align}
For the second estimate, center the identity \(V_n=T_n-1+E_n\) and
apply the reverse triangle inequality in \(L^2\), using
\(\sd(E_n)\le(\EE E_n^2)^{1/2}\). Similarly,
\begin{equation}\label{eq:three-last-cap-bridge}
 0\le V_n-W_n\le6n\ind_{\{H_n>10n\}},\qquad
 0\le\mu_n-\nu_n\le6\rho^n.
\end{equation}
Thus \(\nu_n\to\alpha_3\), since the main proof already gives
\(\mu_n\to\alpha_3\).

\subsection{Numerical forcing bound and concentration}

The replacement parameters of Lemma~\ref{lem:replacement} become
\(v=n-1\), \(\varepsilon=n^{-2/3}\), and
\(g=\lceil6n^{1/3}\rceil\), for \(n\ge4\).
For a satisfiable background \(B\), put \(b=b(B)\) and
\(x=q_2(B)=b(b-1)/(4v(v-1))\). When \(x>\varepsilon\), the proof of
that lemma gives \(b\ge3\), and the exact killing probabilities read
\begin{equation}\label{eq:three-killing}
 q_3(B)=x^{3/2}
 \frac{b-2}{\sqrt{b(b-1)}}
 \frac{\sqrt{v(v-1)}}{v-2}
 \ge\frac13x^{3/2}.
\end{equation}
Here \(b-2\ge b/3\), \(\sqrt{b(b-1)}\le b\), and
\(\sqrt{v(v-1)}\ge v-2\). Thus \(gq_3(B)>2x\), and the lemma's
independent-block estimate is
\(1-(1-q_3(B))^g\ge1-e^{-2x}\ge x\).
Together with the case \(x\le\varepsilon\), this recovers the uniform
replacement loss \(\varepsilon\).

The numerical forcing bound follows directly from
\eqref{eq:forcing-general-cap} with \(k=3\) and \(M=10n\):
\begin{equation}\label{eq:three-forcing}
 \sum_{j=0}^{10n}\EE\frac{b(F_{n,j})}{n}
 \le30e^{30}\lceil6n^{1/3}\rceil
       +(150n+15)n^{-2/3}
 =O(n^{1/3}).
\end{equation}
The constant \(30\) is the maximum mean incidence count \(3M/n\).
The general proof supplies the conditioning and telescoping needed
for this estimate, including the prefix \(j=10n\).

Concentration at the smaller cap follows even more directly from the
main variance bound: \(W_n=\min\{V_n,10n\}\). If \(V_n'\) is an
independent copy of \(V_n\) and \(W_n'=\min\{V_n',10n\}\), truncation
is \(1\)-Lipschitz, so
\[
 \Var(W_n)=\tfrac12\EE(W_n-W_n')^2
 \le\tfrac12\EE(V_n-V_n')^2=\Var(V_n).
\]
Proposition~\ref{prop:concentration} and Chebyshev's inequality give
\begin{equation}\label{eq:three-concentration}
 \Var(W_n)=O(n^{5/3}),\qquad
 \PP\bigl(|W_n-n\nu_n|\ge n^{11/12}\bigr)=O(n^{-1/6}).
\end{equation}
The same powers hold for \(V_n\), with its own finite-size mean.

\subsection{The comparison range is independent of the cap}

In the auxiliary model, an ordered triple has distinct variable indices
with probability
\[
 \vartheta_n=\frac{(n)_3}{n^3}=1-\frac3n+\frac2{n^2}.
\]
At density \(c\), Poisson thinning separates a proper-clause count
\(J\sim\Pois(cn\vartheta_n)\) from an independent repeated-variable
family with mean count \(c(3-2/n)\). Therefore
Lemma~\ref{lem:proper-transfer} gives, for every fixed \(D\ge0\) and
\(0\le c\le D\),
\begin{equation}\label{eq:three-density-sandwich}
 e^{-3D}\EE P_n(J)\le s_n(c)\le\EE P_n(J),\qquad
 cn-3D\le\EE J\le cn,\qquad \Var(J)\le Dn.
\end{equation}
The parameter \(D\) controls the density at which probabilities are
compared. The proof of Theorem~\ref{thm:concentration-to-convergence}
requires \(D>b\); it need not equal the cap.

By \eqref{eq:three-anchors} and \eqref{eq:three-concentration}, both
\((\lambda,x_n^{(\lambda)})=(10,\nu_n)\) and
\((\lambda,x_n^{(\lambda)})=(16,\mu_n)\) satisfy that theorem's inputs
with
\begin{equation}\label{eq:three-engine-parameters}
 (a,b,\delta,\eta)=\left(a_0,7,\frac1{12},\frac16\right),\qquad
 (D,R,\gamma,\kappa)=\left(8,24,\frac16,\frac{e^{-24}}2\right).
\end{equation}
Here \(b<\lambda\) for either cap, \(D>b\), and
\(\gamma=\min\{\eta,1-2\delta\}=1/6\).
On \([0,8]\), the repetition mean is at most \(24\), so the sandwich is
\begin{equation}\label{eq:three-poisson-sandwich}
 e^{-24}\EE P_n(J)\le s_n(c)\le\EE P_n(J),\qquad
 cn-24\le\EE J\le cn,\qquad \Var(J)\le8n.
\end{equation}
For either choice \(z_n=\nu_n\) or \(z_n=\mu_n\), there is an absolute
constant \(K_8<\infty\) such that, uniformly on \(0\le c\le8\) for all
sufficiently large \(n\),
\begin{align}
 c\le z_n-2n^{-1/12}&\Longrightarrow s_n(c)\ge e^{-24}/2,
       \label{eq:three-transfer-low}\\
 c\ge z_n+2n^{-1/12}&\Longrightarrow s_n(c)\le K_8n^{-1/6}.
       \label{eq:three-transfer-high}
\end{align}
The theorem also gives the corresponding size comparisons for either
choice of center:
\begin{align}
 z_{r+s}&\ge\min\{z_r,z_s\}-4\min\{r,s\}^{-1/12},
       \label{eq:three-center-sum}\\
 z_{s+1}&\ge z_s-4s^{-1/12},
       \label{eq:three-neighbor}
\end{align}
for all sufficiently large relevant sizes, without a restriction on
\(r/s\) in the first inequality.

If probabilities are instead compared over the wider interval
\([0,16]\), the same calculation gives \(R=48\), lower sandwich factor
\(e^{-48}\), and lower-transfer constant \(e^{-48}/2\). These are the
constants obtained from the generic choices \(\lambda=16\), \(b=9\),
and \(D=16\) in the main proof. The sharper mean bound permits \(D=8\)
for that same cap, while the \(D=16\) calculation covers a larger range
of densities.
